% Claim proof file for row claim_3_25 -- "ISigmaSeparation".
% Auto-generated stub by scaffold/solving_current_row.py at row start. A
% manager agent dedicated to writing the TeX proof fills in the body below.
% The proof must:
%   - Stay close to the lecture notes' proof if one exists (always search
%     the LN first for a `\begin{proof}` block immediately following the
%     claim; copy / lightly edit when present).
%   - Otherwise use the LN paradigm and cite prior claims by their `ref`.
%   - Be self-contained enough that a mathematician can follow it without
%     re-reading the LN.
%   - Have no `sorry`, `omitted`, or "obvious" shortcuts.
%
% Compiles standalone via the `subfiles` package.

\documentclass[main]{subfiles}

\begin{document}
% ref: claim_3_25 (proof, with statement restated for readability)
% title: ISigmaSeparation
\def\rowref{claim\_3\_25}
\phantomsection\label{claim_3_25}
%
% Cross-references: see claim_<ref>_statement_<title>.tex in this folder
% for the corresponding statement.

% --- statement (restated) ---------------------------------------------------
\begin{Lem}[$i\sigma$-separation under marginalization]\label{lem:stability_separation_marginalization}
        Let $G=(J,V,E,L)$ be a CDMG, $A,B,C \ins J \cup V$ and $D \ins V$ 
        be subsets of nodes such that:
        \[ D \cap \lp A \cup B \cup C\rp = \emptyset.\]
        Then we have the equivalence:
        \[ A \isPerp_G B \given C \qquad \iff \qquad  A \isPerp_{G^{\sm D} }B \given C. \]
\end{Lem}

% --- proof ------------------------------------------------------------------
\begin{proof}
% DISPROVE MODE.  This proof replaces the LN's own proof of
% \texttt{lem:stability\_separation\_marginalization}
% (\texttt{lecture-notes/lecture\_notes/graphs.tex} lines 1424--1578)
% with a proof of the \emph{negation}.  A concrete counter-example was
% constructed by an earlier sub-task worker and independently verified
% by an exhaustive pen-and-paper recomputation; see
% \texttt{workspace\_claim\_3\_25.md} (Sub-task 3 escalation report,
% Sub-task 4 verifier report).  The breaking step in the LN's proof is at
% \texttt{graphs.tex:1559} (= line 172 of the prior version of this proof
% file), which asserts that the fork bifurcation
% $w \hut u \tuh b_{j+1}$ yields a bidirected edge $w \huh b_{j+1} \in
% L^{\sm u}$.  Under the present encoding of marginalization
% (\refrow{def_3_14}, Lean
% \texttt{Section3\_2/Marginalization.lean:530--535}) the
% $L^{\sm W}$ set excludes any pair already supplied by a directed walk
% through $W$, and the directed walk $w \tuh u \tuh b_{j+1}$ in our
% witness forces $(w, b_{j+1}) \in E^{\sm u}$ -- hence
% $(w, b_{j+1}) \notin L^{\sm u}$, and the LN's rerouting step is
% unavailable.  This is the predicted failure mode noted in
% \texttt{claude.md} ("\emph{there does exist at least one claim in the
% lecture notes that is false}").  We disprove the lemma by exhibiting a
% concrete CDMG, sets $A, B, C, D$, and a direct comparison witness.
%
% Cited definitions and prior claims:
%   \refrow{def_3_1}  -- CDMG
%   \refrow{def_3_5}  -- $\Anc^G$, $\Desc^G$, $\Sc^G$
%   \refrow{def_3_14} -- Marginalization $G^{\sm W}$
%   \refrow{def_3_17} -- $\sigma$-blocked / $\sigma$-open walks
%   \refrow{def_3_18} -- $i\sigma$-separation

We disprove the lemma by exhibiting a CDMG $G$ and subsets
$A, B, C \ins J \cup V$, $D \ins V$, satisfying the disjointness
hypothesis $D \cap (A \cup B \cup C) = \emptyset$, but for which the
asserted equivalence \emph{fails}.  Concretely we show
\[ A \nisPerp_G B \given C
   \qquad\text{and}\qquad
   A \isPerp_{G^{\sm D}} B \given C, \]
so the $\Longleftarrow$ direction of the equivalence is broken.

\paragraph{Counter-example setup.}
Let $\alpha$ be any type with at least six distinct elements, label six of
them $v_0, b_j, u, w, b_{j+1}, v_n$, and define the CDMG
$G = (J, V, E, L)$ by:
\begin{align*}
  J &:= \emptyset,\\
  V &:= \lC v_0,\ b_j,\ u,\ w,\ b_{j+1},\ v_n \rC,\\
  E &:= \lC (v_0, b_j),\ (b_j, u),\ (u, b_{j+1}),\ (b_{j+1}, v_n),\
              (u, w),\ (w, u),\ (w, b_j) \rC,\\
  L &:= \emptyset.
\end{align*}

\emph{Verification of the CDMG axioms (\refrow{def_3_1}).}  By inspection,
$J \cap V = \emptyset$ since $J = \emptyset$.  Each pair in $E$ has both
coordinates in $V$, so $E \ins (J \cup V) \times V$.  $L \ins V \times V$
holds vacuously since $L = \emptyset$, as do $L$-irreflexivity, $L$-symmetry,
and $E \cap L = \emptyset$.  Hence $G$ is a CDMG.

Now choose
\[ A := \lC v_0 \rC, \qquad
   B := \lC v_n \rC, \qquad
   C := \lC b_j,\ w \rC, \qquad
   D := \lC u \rC. \]
All four are subsets of $J \cup V = V$, and $D \ins V$.  Furthermore
\[ D \cap (A \cup B \cup C)
   = \lC u \rC \cap \lC v_0,\ v_n,\ b_j,\ w \rC = \emptyset, \]
since $u$ is distinct from each of $v_0, v_n, b_j, w$ by the labelling.
The disjointness hypothesis of the lemma holds.

\paragraph{LHS: $A \nisPerp_G B \given C$.}
We exhibit a $C$-$\sigma$-open walk in $G$ from $v_0 \in A$ to
$v_n \in J \cup B = \lC v_n \rC$.  Take
\[ \pi := \lp v_0 \tuh b_j \tuh u \tuh b_{j+1} \tuh v_n \rp, \]
a walk in $G$ of length $4$, since each of the four ordered pairs
$(v_0, b_j),\ (b_j, u),\ (u, b_{j+1}),\ (b_{j+1}, v_n)$ lies in $E$.

\emph{Auxiliary computations of $\Anc^G$, $\Desc^G$, $\Sc^G$
(\refrow{def_3_5}).}  By definition $\Anc^G(v) = \lC u \in J \cup V :
\text{there is a directed walk } u \tuh \cdots \tuh v \text{ in } G \rC$
(including $v$ itself via the length-$0$ walk), and dually for
$\Desc^G$; the strongly connected component is
$\Sc^G(v) = \Anc^G(v) \cap \Desc^G(v)$.  Tracing edges of $G$:
\begin{itemize}
\item $\Anc^G(b_j) = \lC v_0,\ b_j,\ u,\ w \rC$: witnesses are
  $v_0 \tuh b_j$ (direct), $w \tuh b_j$ (direct), and $u \tuh w \tuh b_j$.
\item $\Desc^G(b_j) = \lC b_j,\ u,\ w,\ b_{j+1},\ v_n \rC$: witnesses are
  $b_j \tuh u$, $b_j \tuh u \tuh w$, $b_j \tuh u \tuh b_{j+1}$, and
  $b_j \tuh u \tuh b_{j+1} \tuh v_n$.
\item Hence $\Sc^G(b_j) = \lC b_j,\ u,\ w \rC$.
\item $\Anc^G(u) = \lC v_0,\ b_j,\ u,\ w \rC$ ($w \tuh u$ direct;
  $b_j \tuh u$ direct; $v_0 \tuh b_j \tuh u$);
  $\Desc^G(u) = \lC u,\ w,\ b_j,\ b_{j+1},\ v_n \rC$ ($u \tuh w$,
  $u \tuh w \tuh b_j$, $u \tuh b_{j+1}$, $u \tuh b_{j+1} \tuh v_n$).
  Hence $\Sc^G(u) = \lC u,\ b_j,\ w \rC$.
\item For $b_{j+1}$ the only outgoing edge is $(b_{j+1}, v_n)$ and $v_n$
  has no outgoing edges, so $\Desc^G(b_{j+1}) = \lC b_{j+1},\ v_n \rC$.
  The only predecessor of $b_{j+1}$ in $E$ is $u$, so working backwards
  $\Anc^G(b_{j+1}) = \lC v_0,\ b_j,\ u,\ w,\ b_{j+1} \rC$.  Hence
  $\Sc^G(b_{j+1}) = \lC b_{j+1} \rC$.
\end{itemize}

\emph{Per-position $\sigma$-open verification on $\pi$
(\refrow{def_3_17}).}  Every step of $\pi$ is forward; a forward step
$a \tuh b$ contributes an arrowhead at $b$ (the target) and \emph{no}
arrowhead at $a$ (the source).  The endpoints $v_0$ and $v_n$ are not in
$C$ (since $C = \lC b_j, w \rC$ contains neither), discharging the
endpoint clauses of $\sigma$-open in the standard sense.

\begin{itemize}
\item \textbf{Position 1 ($b_j$).}  Step-in $v_0 \tuh b_j$ (arrowhead at
  $b_j$), step-out $b_j \tuh u$ (no arrowhead at $b_j$ from this side).
  Only one arrowhead meets $b_j$, so this is a \emph{non-collider}; the
  local pattern matches "right chain" $v_0 \suh b_j \tuh u$ in def_3_17.
  Unblockability for a right chain at $b_j$ requires the forward-out
  target to lie in $\Sc^G(b_j)$; here the target is $u$, and
  $u \in \Sc^G(b_j) = \lC b_j, u, w \rC$.  Therefore $b_j$ is an
  \emph{unblockable} non-collider on $\pi$.  By def_3_17 clause i.ii
  (only blockable non-colliders in $C$ obstruct $\sigma$-open),
  $b_j \in C$ does \emph{not} disqualify this position.
\item \textbf{Position 2 ($u$).}  Step-in $b_j \tuh u$, step-out
  $u \tuh b_{j+1}$.  Non-collider (right chain $b_j \tuh u \tuh b_{j+1}$),
  forward-out target $b_{j+1} \notin \Sc^G(u) = \lC u, b_j, w \rC$;
  hence a \emph{blockable} non-collider.  Since $u \notin C$, clause
  i.ii is satisfied at this position.
\item \textbf{Position 3 ($b_{j+1}$).}  Step-in $u \tuh b_{j+1}$, step-out
  $b_{j+1} \tuh v_n$.  Non-collider (right chain), forward-out target
  $v_n \notin \Sc^G(b_{j+1}) = \lC b_{j+1} \rC$; blockable.  Since
  $b_{j+1} \notin C$, clause i.ii is satisfied.
\end{itemize}

\emph{No colliders on $\pi$.}  A collider at position $k$ requires two
arrowheads at $v_k$, one from the step-in (its target side) and one from
the step-out (its source side).  Since every step of $\pi$ is forward,
each step-out contributes no source-side arrowhead.  Thus no interior
position of $\pi$ is a collider, and def_3_17 clause i.i is vacuous.

Combining the three interior positions and the absence of colliders,
$\pi$ is $C$-$\sigma$-open.  By \refrow{def_3_18} (negation), the
existence of this walk gives $A \nisPerp_G B \given C$.

\paragraph{RHS: $A \isPerp_{G^{\sm u}} B \given C$.}
We must show that \emph{every} walk in $G^{\sm u}$ from $v_0 \in A$ to
$v_n \in J^{\sm u} \cup B = \lC v_n \rC$ is $C$-$\sigma$-blocked.  We
first compute the components of $G^{\sm u}$ explicitly.

\subparagraph{Computing $E^{\sm u}$.}  By \refrow{def_3_14},
$(\ul v, \ol v) \in E^{\sm u}$ iff $\ul v, \ol v \in J \cup V \sm \lC u \rC$
and there exists a directed walk in $G$
\[ \ul v \tuh w_1 \tuh \cdots \tuh w_{n-1} \tuh \ol v \]
with every intermediate $w_i \in \lC u \rC$.  We enumerate:
\begin{itemize}
\item \emph{Length $1$} (no intermediate): the original edges of $E$ whose
  endpoints both lie in $V \sm \lC u \rC$, namely $(v_0, b_j),\
  (b_{j+1}, v_n),\ (w, b_j)$.
\item \emph{Length $2$ via $u$}: pairs $(x, y)$ with $(x, u),\ (u, y) \in E$
  and $x, y \in V \sm \lC u \rC$.  Predecessors of $u$ in $E$: $\lC b_j,
  w \rC$.  Successors of $u$ in $E$: $\lC b_{j+1}, w \rC$.  The
  Cartesian product yields $(b_j, b_{j+1}),\ (b_j, w),\ (w, b_{j+1}),\
  (w, w)$.
\item \emph{Length $\ge 3$ via $u$}: a directed walk of length $\ge 3$
  with interior in $\lC u \rC$ would have to repeat $u$, requiring an
  edge $(u, u) \in E$.  No such edge exists in our $E$.  So no length
  $\ge 3$ walks contribute.
\end{itemize}
Putting it together,
\[
  E^{\sm u}
  = \lC (v_0, b_j),\ (b_{j+1}, v_n),\ (w, b_j),\
         (b_j, b_{j+1}),\ (b_j, w),\ (w, b_{j+1}),\
         (w, w) \rC.
\]

\subparagraph{Computing $L^{\sm u}$: the load-bearing step.}  In the
Lean formalization of marginalization
(\refrow{def_3_14}; Lean
\texttt{Section3\_2/Marginalization.lean:530--535}) the set $L^{\sm u}$
consists of those pairs $(\ul v, \ol v)$ with
$\ul v, \ol v \in V \sm \lC u \rC$, $\ul v \ne \ol v$, satisfying
\emph{both} of the exclusion clauses
\begin{align*}
  &\neg \exists \pi : \text{Walk } G\ \ul v\ \ol v,
       \quad \pi \text{ is directed with interior in } \lC u \rC, \\
  &\neg \exists \pi : \text{Walk } G\ \ol v\ \ul v,
       \quad \pi \text{ is directed with interior in } \lC u \rC,
\end{align*}
and \emph{some} bifurcation between $\ul v$ and $\ol v$ in $G$ with
interior in $\lC u \rC$ (in either walk direction).  The two exclusion
clauses are forced by \refrow{def_3_1}'s $\texttt{disjoint\_EL}$ field
($E \cap L = \emptyset$), which prevents a pair already supplied by a
directed walk through $u$ from also being recorded as a bidirected edge
of $G^{\sm u}$.\footnote{This is the encoding-level deviation from the
LN's literal definition of $G^{\sm W}$ in \refrow{def_3_14}, openly documented
in the design block at
\texttt{Section3\_2/Marginalization.lean:373--425}.  The LN's framework
treats $E$ and $L$ as different sorts of objects, so a pair can be both
a directed edge and a bidirected edge of $G^{\sm u}$; under the Lean
encoding $E^{\sm u}$ subsumes the directed-walk contribution and
$L^{\sm u}$ records only the bidirected contributions not already
covered by $E^{\sm u}$.}

We claim $L^{\sm u} = \emptyset$.  We enumerate, then exclude, every
candidate.

\emph{Enumeration of candidate pairs.}  Since $G$ has $L = \emptyset$, a
bifurcation in $G$ with interior in $\lC u \rC$ uses only directed edges,
read forward or backward.  The relevant bifurcation shapes (in the
generalised sense: a bifurcation walk is one whose steps split into a
"left arm" of zero or more backward (and optionally a single bidirected)
steps and a "right arm" of zero or more forward steps) are:
\begin{itemize}
\item \emph{Length $1$, single backward step} $\ol v \hut \ul v$ from
  $(\ul v, \ol v) \in E$ read backwards.  Candidate unordered pairs:
  $\lC v_0,\ b_j \rC$ from $(v_0, b_j)$;
  $\lC b_{j+1},\ v_n \rC$ from $(b_{j+1}, v_n)$;
  $\lC w,\ b_j \rC$ from $(w, b_j)$.
\item \emph{Length $2$, double-backward} $\ul v \hut u \hut \ol v$,
  requiring $(u, \ul v),\ (u, \ol v) \in E$.\footnote{An all-backward
  walk is a degenerate bifurcation with empty right arm; see Lean
  \texttt{Section3\_1/Bifurcation.lean} and the discussion at
  \texttt{Section3\_2/Marginalization.lean:401--407}.}  Successors of
  $u$ in $E$: $\lC b_{j+1},\ w \rC$, so the unordered candidates are
  $\lC b_{j+1},\ w \rC$ (and the trivial $\lC b_{j+1} \rC$,
  $\lC w \rC$ which violate $\ul v \ne \ol v$).
\item \emph{Length $2$, fork} $\ul v \hut u \tuh \ol v$, requiring
  $(u, \ul v),\ (u, \ol v) \in E$.  Same successors as above; candidate
  unordered pair $\lC b_{j+1},\ w \rC$ again (the asymmetric pairs
  $(b_{j+1}, w)$ and $(w, b_{j+1})$ collapse to one unordered pair).
\item \emph{Length $2$, hinge involving an $L$-edge}: impossible since
  $L = \emptyset$.
\item \emph{Length $2$, double-forward} $\ul v \tuh u \tuh \ol v$: this
  is a directed walk, not a bifurcation (the hinge step would have to
  carry an arrowhead at its source, which a forward step does not).  So
  it contributes only to $E^{\sm u}$, not to $L^{\sm u}$.
\item \emph{Length $\ge 3$}: as in the $E^{\sm u}$ enumeration, such a
  walk would need to repeat $u$ in its interior, which requires
  $(u, u) \in E$.  None exists.
\end{itemize}

Collecting unordered candidates we have
\[ \lC v_0, b_j \rC,\ \lC b_{j+1}, v_n \rC,\ \lC w, b_j \rC,\
   \lC b_{j+1}, w \rC. \]

\emph{Exclusion of each candidate via the directed-walk clauses.}  For
each unordered candidate, we exhibit a directed walk in $G$ between the
two endpoints (in some direction) with interior in $\lC u \rC$, which
triggers one of the two exclusion clauses of $L^{\sm u}$:
\begin{itemize}
\item $\lC v_0, b_j \rC$: directed walk $v_0 \tuh b_j$ of length $1$,
  interior empty, witnesses $(v_0, b_j) \in E^{\sm u}$ and excludes the
  $\ul v \to \ol v$ direction.
\item $\lC b_{j+1}, v_n \rC$: directed walk $b_{j+1} \tuh v_n$
  excludes $(b_{j+1}, v_n)$ direction.
\item $\lC w, b_j \rC$: directed walk $w \tuh b_j$ excludes $(w, b_j)$
  direction.
\item $\lC b_{j+1}, w \rC$: the directed walk $w \tuh u \tuh b_{j+1}$
  in $G$ has length $2$ with interior $\lC u \rC$; it witnesses
  $(w, b_{j+1}) \in E^{\sm u}$ and excludes the $(w, b_{j+1})$ direction
  of the $L^{\sm u}$ clause.
\end{itemize}

Every candidate pair is excluded; hence $L^{\sm u} = \emptyset$.

\subparagraph{Computing $\Sc^{G^{\sm u}}$ and $\Anc^{G^{\sm u}}(C)$.}
We now have $J^{\sm u} = \emptyset$,
$V^{\sm u} = \lC v_0,\ b_j,\ w,\ b_{j+1},\ v_n \rC$,
$E^{\sm u}$ as enumerated, and $L^{\sm u} = \emptyset$.  Ancestral and
descendant sets in $G^{\sm u}$ are computed from $E^{\sm u}$ alone (no
bidirected edges interact with directedness for $\Anc/\Desc$):
\begin{itemize}
\item $\Desc^{G^{\sm u}}(b_j)$: starting from $b_j$ we can take
  $(b_j, b_{j+1}),\ (b_j, w)$; from $w$ also $(w, b_j),\ (w, w),\
  (w, b_{j+1})$; from $b_{j+1}$ also $(b_{j+1}, v_n)$.  So
  $\Desc^{G^{\sm u}}(b_j) = \lC b_j,\ w,\ b_{j+1},\ v_n \rC$.
\item $\Anc^{G^{\sm u}}(b_j)$: predecessors of $b_j$ in $E^{\sm u}$ are
  $v_0$ and $w$; predecessors of $w$ are $b_j$ and $w$.  Hence
  $\Anc^{G^{\sm u}}(b_j) = \lC v_0,\ b_j,\ w \rC$.
\item $\Sc^{G^{\sm u}}(b_j) =
       \lC v_0, b_j, w \rC \cap \lC b_j, w, b_{j+1}, v_n \rC
     = \lC b_j,\ w \rC$.
\item By symmetric calculation,
  $\Anc^{G^{\sm u}}(w) = \lC v_0,\ b_j,\ w \rC$,
  $\Desc^{G^{\sm u}}(w) = \lC w,\ b_j,\ b_{j+1},\ v_n \rC$, and
  $\Sc^{G^{\sm u}}(w) = \lC b_j,\ w \rC$.
\item $\Desc^{G^{\sm u}}(b_{j+1}) = \lC b_{j+1},\ v_n \rC$ (only
  $(b_{j+1}, v_n)$ leaves $b_{j+1}$).  Predecessors of $b_{j+1}$ in
  $E^{\sm u}$ are $b_j$ and $w$, hence
  $\Anc^{G^{\sm u}}(b_{j+1}) = \lC v_0,\ b_j,\ w,\ b_{j+1} \rC$.
  So $\Sc^{G^{\sm u}}(b_{j+1}) = \lC b_{j+1} \rC$.
\item $\Desc^{G^{\sm u}}(v_n) = \lC v_n \rC$ (no outgoing edges from
  $v_n$ in $E^{\sm u}$).
\end{itemize}
Therefore
\[ \Anc^{G^{\sm u}}(C)
   = \Anc^{G^{\sm u}}(b_j) \cup \Anc^{G^{\sm u}}(w)
   = \lC v_0,\ b_j,\ w \rC. \]
The salient consequences are
\[ b_{j+1} \notin \Sc^{G^{\sm u}}(b_j),
   \quad b_{j+1} \notin \Sc^{G^{\sm u}}(w),
   \quad v_n \notin \Anc^{G^{\sm u}}(C). \]

\subparagraph{Exhaustive walk-by-walk blocking.}
Let $\pi' = \lp w_0 \sus w_1 \sus \cdots \sus w_n \rp$ be \emph{any} walk
in $G^{\sm u}$ with $w_0 = v_0$ and $w_n = v_n$; necessarily $n \ge 1$.
Since $L^{\sm u} = \emptyset$, every step of $\pi'$ is either forward
(via $E^{\sm u}$) or backward (via $E^{\sm u}$ read in reverse).  No
bidirected steps occur.

\emph{Forced final two steps.}  Inspecting $E^{\sm u}$, the only edge of
$G^{\sm u}$ that touches $v_n$ is $(b_{j+1}, v_n)$.  So $v_n$ has exactly
one predecessor ($b_{j+1}$) and exactly one successor ($b_{j+1}$, via
the same edge read backward).  Hence the final step
$w_{n-1} \sus w_n = w_{n-1} \sus v_n$ is one of:
\begin{itemize}
\item forward $w_{n-1} \tuh v_n$ with $w_{n-1} = b_{j+1}$, via
  $(b_{j+1}, v_n) \in E^{\sm u}$; or
\item backward $w_{n-1} \hut v_n$ -- impossible, since this would
  require some edge $(v_n, w_{n-1}) \in E^{\sm u}$, but no edge of
  $E^{\sm u}$ has $v_n$ as its first coordinate.
\end{itemize}
The final step is therefore forced to be forward $b_{j+1} \tuh v_n$, and
$w_{n-1} = b_{j+1}$.

\emph{Penultimate step into $b_{j+1}$.}  Now case-split on the step
$w_{n-2} \sus w_{n-1} = w_{n-2} \sus b_{j+1}$.  Predecessors of
$b_{j+1}$ in $E^{\sm u}$ are exactly $b_j$ and $w$ (from
$(b_j, b_{j+1}),\ (w, b_{j+1})$).  Successors of $b_{j+1}$ in
$E^{\sm u}$ are exactly $v_n$ (from $(b_{j+1}, v_n)$).  So:
\begin{itemize}
\item[(A)] forward $w_{n-2} \tuh b_{j+1}$ with $w_{n-2} \in \lC b_j,\ w \rC$;
\item[(B)] backward $w_{n-2} \hut b_{j+1}$, i.e.\ $(b_{j+1}, w_{n-2})
  \in E^{\sm u}$, forces $w_{n-2} = v_n$.
\end{itemize}
(Case C, a bidirected step, is excluded by $L^{\sm u} = \emptyset$.)

\textbf{Case (A): $w_{n-2} \in \lC b_j, w \rC$.}
At position $n-2$ of $\pi'$ the step-out is forward
$w_{n-2} \tuh b_{j+1}$, contributing \emph{no} source-side arrowhead at
$w_{n-2}$.  Whatever the step-in $w_{n-3} \sus w_{n-2}$ may be, the
collider condition (an arrowhead from each side meeting at $w_{n-2}$)
fails, so position $n-2$ is a \emph{non-collider}.  For unblockability
of a non-collider whose forward-out target is $b_{j+1}$, def_3_17
requires $b_{j+1} \in \Sc^{G^{\sm u}}(w_{n-2})$ (right-chain clause).
However,
$\Sc^{G^{\sm u}}(b_j) = \Sc^{G^{\sm u}}(w) = \lC b_j,\ w \rC$, and
$b_{j+1} \notin \lC b_j,\ w \rC$.  Hence position $n-2$ is a
\emph{blockable} non-collider.  Now $w_{n-2} \in \lC b_j, w \rC = C$, so
def_3_17 clause ii.ii fires: $\pi'$ contains a blockable non-collider in
$C$ and is $C$-$\sigma$-blocked.

\textbf{Case (B): $w_{n-2} = v_n$.}
The step-out from position $n-2$ is backward $v_n \hut b_{j+1}$, which
\emph{does} contribute a source-side arrowhead at $v_n$.  The step-in
$w_{n-3} \sus w_{n-2} = w_{n-3} \sus v_n$: as noted in the
"forced-final-two-steps" analysis, the only edge of $E^{\sm u} \cup
L^{\sm u}$ touching $v_n$ is $(b_{j+1}, v_n)$, hence $w_{n-3} = b_{j+1}$
and the step-in is forward $b_{j+1} \tuh v_n$, contributing a
target-side arrowhead at $v_n$.  Both step-in and step-out have
arrowheads at $v_n$, so $v_n$ is a \emph{collider} on $\pi'$.  For
$\sigma$-open at a collider, def_3_17 requires $v_n \in \Anc^{G^{\sm u}}(C)$.
But $\Anc^{G^{\sm u}}(C) = \lC v_0,\ b_j,\ w \rC$ does not contain
$v_n$.  Hence def_3_17 clause ii.i fires: $\pi'$ contains a collider
outside $\Anc^{G^{\sm u}}(C)$ and is $C$-$\sigma$-blocked.

\emph{Conclusion of the case analysis.}  In each of the two exhaustive
cases (A) and (B), $\pi'$ is $C$-$\sigma$-blocked.  Since the analysis
depends only on the final two steps of $\pi'$, this holds independently
of the walk's length or earlier structure -- arbitrary detours through
the $(w, w)$ self-loop, back-and-forth through $b_j$ and $w$, or
revisits to $v_n$ all eventually leave $\pi'$ in one of cases (A) or
(B) on its final approach to $v_n$, and the blocking position arises
there.

Hence \emph{every} walk in $G^{\sm u}$ from $v_0$ to $v_n$ is
$C$-$\sigma$-blocked, and by \refrow{def_3_18}
$A \isPerp_{G^{\sm u}} B \given C$.

\paragraph{Conclusion.}
With $D = \lC u \rC$ we have shown
$A \nisPerp_G B \given C$ and $A \isPerp_{G^{\sm D}} B \given C$.  Hence
the equivalence
\[ A \isPerp_G B \given C \quad\iff\quad A \isPerp_{G^{\sm D}} B \given C \]
fails on this witness; specifically, the $\Longleftarrow$ direction is
broken.  The lemma as stated is therefore false.  \qed

\paragraph{Where the LN proof fails (audit trail).}
For completeness, we point out the precise step in the LN's own proof
that breaks on this witness.  In the LN proof
(\texttt{lecture-notes/lecture\_notes/graphs.tex} lines 1424--1578;
= lines 36--190 of the prior version of this file), the $(\impliedby)$
direction constructs $\pi'$ by contracting each maximal $u$-run on a
$C$-$\sigma$-open walk in $G$ and then rerouting at unblockable
non-colliders that become blockable in $G^{\sm u}$.  The rerouting at
an unblockable non-collider $b_j \in C$ uses a successor $w$ of $u$ on a
directed path $u \tuh \cdots \tuh b_j$ in $G$, and asserts at
\texttt{graphs.tex:1559} (line 172 of the prior proof):
\begin{quote}
\emph{"The fork bifurcation $w \hut u \tuh b_{j+1}$ yields $w \huh
b_{j+1} \in L^{\sm u}$."}
\end{quote}
On our witness this assertion is \emph{false}.  The fork bifurcation
$w \hut u \tuh b_{j+1}$ does exist in $G$ (since $(u, w),\ (u, b_{j+1})
\in E$), but the candidate pair $\lC w,\ b_{j+1} \rC$ is excluded from
$L^{\sm u}$ by the first (or second, by symmetry) exclusion clause of
\refrow{def_3_14}: the directed walk $w \tuh u \tuh b_{j+1}$ in $G$
(length $2$, interior $\lC u \rC$) is a witness for
$(w, b_{j+1}) \in E^{\sm u}$, and the $\texttt{disjoint\_EL}$ constraint
in \refrow{def_3_1} consequently forces $(w, b_{j+1}) \notin L^{\sm u}$.

The LN's rerouted walk
$\lp \cdots \tuh b_j \tuh w \huh b_{j+1} \tuh v_n \rp$
therefore cannot be assembled in $G^{\sm u}$.  Any alternative attempt
to reach $b_{j+1}$ via $w$ must use the forward edge
$(w, b_{j+1}) \in E^{\sm u}$ -- but then position $w$ on $\pi'$ has a
forward-out to $b_{j+1} \notin \Sc^{G^{\sm u}}(w) = \lC b_j,\ w \rC$,
making $w$ a blockable non-collider in $C$ and triggering Case (A) of
the analysis above.

Thus the LN's lemma is false on the present encoding, the broken step
is the bidirected-edge-inclusion claim at \texttt{graphs.tex:1559}, and
the structural reason is the $\texttt{disjoint\_EL}$-driven exclusion
clauses of $L^{\sm u}$.
\end{proof}

\end{document}
